\revision[R2.D9]{
\subsection{Using FDP Placement Hints}\label{sec:fdp}
\module{FDP exposes RU handles for host-side placement hints.}
When an SSD supports FDP, the NoWA pattern becomes unnecessary to achieve an SSD WAF of 1.
In addition to the RU size, FDP-enabled SSDs expose \emph{Reclaim Unit Handles (RUHs)}, each managing a disjoint set of RUs~\cite{fdp}.
The number of RUHs determines how many independent write streams the SSD can maintain without multiplexing.
The host can leverage this by providing placement hints when issuing writes, preventing pages with different deathtimes from being mixed internally~\cite{song_fast26_warp}.
In our design, this directly corresponds to the maximum number of active zones.

\module{Optimization: using FDP to avoid multiplexing.}
FDP placement hints are straightforward to use in our design.
We assign each zone a placement ID (PlID) modulo the number of RUHs and persist this mapping per zone.
For example, with four RUHs (PlIDs 0--3), zones A--D map to PlIDs 0--3, so writes to zone A always use PlID~0.
Under the write pattern shown in \cref{fig:multipleopen}, writes to different zones are directed to different RUHs, preventing interleaving between zones such as A and C (\cref{fig:fdp}).
With the internal SSD information exposed by FDP, an SSD WAF of~1 is achievable as long as the zone size equals the RU size and the number of open zones does not exceed the RUH count (as shown in \cref{tab:fdpresults}). }
